\documentclass[subfile]{~/PhyAI/phyai.tex}
\begin{document}
\chapter{Boltzmann Statistics}\label{chap:boltzmann_statistics} % (fold)
\section{Boltzmann Distribution}\label{sec:boltzmann_distribution} % (fold)
In this chapter, we will learn how to apply some rules of probability and statistics to calculate the probability that a particle is in any given microstate. For example, inside a gas container, we can calculate the probability of finding particles moving with particular speed range, or any other states. Firstly, let's derive the probability density function:
\begin{figure}[H]
	\begin{center}
		\includegraphics[width=0.6\textwidth]{~/PhyAI/chapter5/5_1.png}
	\end{center}
	\caption{An imaginary experiment}\label{fig:5_1}
\end{figure}
Consider a system divided into two parts. The first part is the single particle found in a microstate $s$, and the second part is the rest of the container (or what I call the reservoir or environment). For any microstates $s_{1}$ and $s_{2}$, we have:
\begin{equation*}
	\begin{split}
		\frac{\mathcal{P}(s_{1})}{\mathcal{P}(s_{2})}&=\frac{\Omega_{R}(s_{1})}{\Omega_{R}(s_{2})}=\exp{\left(\frac{S_{R}(s_{1})-S_{R}(s_{2})}{k_{B}}\right)}
	\end{split}
\end{equation*}
From the thermodynamics identity (Theorem ~\ref{The Thermodynamics Identity}):
$$dS=\frac{dU }{T}\quad\text{($V$ and $N$ are held as constants)}$$
For abbreviation, I denote: $$ \beta=\frac{1 }{k_{B}T} $$
We obtain:
\begin{equation*}
	\begin{split}
		\frac{\mathcal{P}(s_{1})}{\mathcal{P}(s_{2})}&=\exp{(\beta(U_{R}(s_{1})-U_{R}(s_{2})))}\\
		&=\exp(-\beta(E(s_{1})-E(s_{2})))\\
		\Leftrightarrow \frac{\mathcal{P}(s_{1}) }{\mathcal{P}(s_{2})}&=\frac{\exp{(-\beta E(s_{1}))}}{\exp(-\beta E(s_{2}))}
	\end{split}
\end{equation*}
Finally for any microstate $s$, we always have the proportional relationship:
$$ \mathcal{P}(s)\propto e^{-\beta E(s)} $$
The probability density $\mathcal{P}(s)$ is formed by normalizing function $Z$ as follows:
$$ Z=\sum_{s}e^{-\beta E(s)}\Rightarrow \mathcal{P}(s)=\frac{1 }{Z}e^{-\beta E(s)} $$
$Z$ is the \textbf{partition function}, and exponential factor $e^{-\beta E(s)}$ is called the \textbf{Boltzmann factor}.
\begin{definition}[Boltzmann Distribution]
	\label{Boltzmann Distribution}
	The probability of finding particle in any microstate $s$ is:
	$$ \mathcal{P}(s)=\frac{1 }{Z}e^{-\beta E(s)} $$
	with $Z$ is the \textbf{partition function}: $$ Z=\sum_{s}e^{-\beta E(s)} $$
\end{definition}
\begin{corollary}[Average Energy]
	\label{Average Energy}
	The average energy of any system is:
	$$ \overline{E}=-\frac{\partial \ln{Z} }{\partial \beta} $$
\end{corollary}
\begin{proof}
	\begin{equation*}
		\begin{split}
			\overline{E}&=\sum_{s}E(s)\mathcal{P}(s)=\frac{1 }{Z}\sum_{s}E(s)e^{-\beta E(s)}=-\frac{1 }{Z}\frac{\partial Z }{\partial \beta}=-\frac{\partial \ln{Z} }{\partial \beta}
		\end{split}
	\end{equation*}
\end{proof}
This above result can be used to derive the equipartition theorem (Theorem ~\ref{The Equipartition Theorem}); for any quadratic degree of freedom, its energy has the form:
$$ E(q)=cq^{2} $$
The partition function $Z$ is:
\begin{equation*}
	\begin{split}
		Z&=\sum_{s}e^{-\beta E(s)}=\sum_{s}e^{-\beta cq^{2}}=\frac{1 }{\Delta q}\int_{-\infty }^{+\infty}e^{-\beta cq^{2}}dq=\frac{1 }{\Delta q\sqrt{\beta c}}\int_{-\infty }^{+\infty}e^{-q^{2}}dq\\
		&=\frac{1 }{\Delta q}\sqrt{\frac{\pi }{\beta c}}
	\end{split}
\end{equation*}
We can see that:
\begin{equation*}
	\begin{split}
		\overline{E}&=-\frac{1 }{Z}\frac{\partial Z }{\partial \beta}=\frac{1 }{2}\beta^{-1}=\frac{1 }{2}k_{B}T
	\end{split}
\end{equation*}
Our obtained result once again, verify the correctness of equipartition theorem.
% section Boltzmann Distribution (end)
\section{Maxwell Speed Distribution}\label{sec:maxwell_distribution} % (fold)
What is the probability of finding particle moving at speed range $(v_{1},v_{2})$? Well, the first thing we need to do is deriving a probability density function of speed $\mathcal{D}(v)$. But here is the subtle observation, a particle moving at speed $v$ has an infinite velocity vector $\vec{v}$ representing how it moves in 3D space. All these velocity vectors form a "velocity sphere", which is similar to the concept of momentum space (as we discussed in Section ~\ref{sec:the_ideal_gas_model}):
\begin{figure}[H]
	\begin{center}
		\includegraphics[width=0.65\textwidth]{~/PhyAI/chapter5/5_2.png}
	\end{center}
	\caption{The velocity sphere and its volume}\label{fig:5_2}
\end{figure}
The probability of finding any particle moving at speed $\vec{v}$ is:
$$ P(\vec{v})\propto \exp{\left(-\beta.\frac{1 }{2}mv^{2}\right)} $$
We are ready to plug every piece and build the $\mathcal{D}(v)$ function:
$$ \mathcal{D}(v)\propto \exp{\left(-\beta.\frac{1}{2}mv^{2}\right)} .4\pi v^{2}$$
The probability density function $\mathcal{D}(v)$ is rewritten as:
$$ \mathcal{D}(v)=C.\exp{\left(-\beta.\frac{1}{2}mv^{2}\right)}.4\pi v^{2} $$
Our final step is determining the constant $C$:
\begin{equation*}
	\begin{split}
		1&=C.\int_{0 }^{+\infty} \exp{\left(-\beta.\frac{1}{2}mv^{2}\right)}.4\pi v^{2}dv\\
		\Rightarrow \frac{1}{C}&=\int_{0 }^{+\infty}\exp{\left(-\beta .\frac{1}{2}mv^{2}\right)}4\pi v^{2}dv \\
		&=4\pi\left(\beta .\frac{1}{2}m\right)^{-3/2}
		\int_{0 }^{+\infty}\exp{\left[-\left(v \sqrt{\beta.\frac{1}{2}m}\right)^{2}\right]} \left(v \sqrt{\beta .\frac{1}{2}m}\right)^{2}.d \left(v\sqrt{\beta.\frac{1}{2}m}\right)\\
		&=4\pi \left(\beta.\frac{1 }{2}m\right)^{-3/2}\int_{0}^{+\infty}e^{-v^{2}}v^{2}dv
	\end{split}
\end{equation*}
Integrating by part yields:
\begin{equation*}
	\begin{split}
		\int_{0}^{+\infty}e^{-v^2}v^{2}dv&=-\frac{1 }{2}.\int_{0 }^{+\infty}v.d(e^{-v^2})=\frac{1 }{2}\int_{0 }^{+\infty}e^{-v^{2}}dv=\frac{\sqrt{\pi }}{4}
	\end{split}
\end{equation*}
We obtain: $$ C=\left(\frac{\beta .m}{2\pi}\right)^{3/2} $$
\begin{definition}[Maxwell Speed Distribution]
	The probability density function of finding any particle at a particular speed range is:
	$$ \mathcal{D}(v)= \left(\frac{\beta.m}{2\pi}\right)^{3/2}.\exp{\left(\frac{-\beta mv^{2}}{2}\right)}.4\pi v^{2}$$
\end{definition}
% section Maxwell Distribution (end)


% section Partition Function and Helmholtz Free Energy (end)
% chapter Boltzmann Statistics (end) 
\end{document}
